\documentclass[11pt]{article}
\usepackage{fontenc}
\usepackage{fontspec}
\setmainfont{Latin Modern Roman}

\begin{document}

Texto Unicode compilado con XeLaTeX: áéíóú ñ ü ¿qué? ¡hola! «comillas latinas» “comillas inglesas” — raya — y símbolos como € § ¶ ° ½.

Un párrafo largo en una sola línea para que el sync y el soft wrap tengan algo que hacer también con este motor: XeTeX genera un \texttt{.xdv} intermedio y el \texttt{.synctex.gz} lo deja xdvipdfmx, de modo que las coordenadas pueden diferir de las de pdfTeX y conviene medirlas aparte en las pruebas de SyncTeX.

\end{document}
